\chapter{Introduction}

\section{Problem Statement}
Autonomous drone systems are increasingly deployed in applications such as infrastructure inspection \cite{jacobsen2023drone},
surveillance \cite{dinesh2018autonomous}, and environmental monitoring \cite{anderson2013uav}. While reliable low-level flight 
control and sensing capabilities have been widely developed, expressing high-level mission objectives in a
way that is both intuitive for human operators and executable by autonomous systems remains
a challenge.

Existing approaches rely on predefined workflows or formal planning models, which require expert knowledge and lack flexibility
when faced with dynamic, changing environments. As a result, translating human instructions into concrete mission plans
remains a bottleneck.

Recent advances in Large Language Models (LLMs) offer a promising approach for bridging
this gap. LLMs can interpret natural-language instructions and infer task structure, enabling
them to transform human mission descriptions into structured plans that autonomous systems
can execute. This capability makes them well suited for high-level reasoning tasks such as task
decomposition, allocation, and scheduling in multi-agent systems.
This project proposes the use of LLMs to generate executable plans for multi-drone missions.
Given a natural language mission description, the system decomposes the task into atomic
subtasks, allocates them to drones according to their capabilities, and generates a structured
execution schedule.

To address the challenges of real-world deployment, the system adopts a hierarchical architecture
in which cloud-based LLMs perform high-level mission planning, while smaller locally deployed models
monitor system state and enable adaptive replanning in response to failures, delays, or environmental changes.
The goal is to develop a flexible and robust framework capable of translating human intent into coordinated 
multi-drone behavior.

The implementation is publicly available to support reproducibility and further development.
\footnote{The source code developed for this thesis is available at: 
\url{https://github.com/Akos-Varga/llm-drone-planning-replanning}.}

\section{Research Goals}

The primary goal of this thesis is to investigate the feasibility and effectiveness of using
Large Language Models for high-level planning and adaptive replanning in multi-drone systems.
While LLMs have demonstrated strong capabilities in natural language understanding and reasoning,
their applicability to real-time, safety-critical autonomous systems remains an open question.

This thesis aims to explore the following research goals:

\begin{itemize}
    \item To evaluate how effectively LLMs can translate natural language mission descriptions into structured,
    executable task plans for multi-drone systems.
    
    \item To investigate the ability of LLMs to perform task decomposition, allocation, and scheduling in environments
    with heterogeneous agents and constrained resources.
    
    \item To assess the robustness of LLM-based planning when operating in dynamic environments that require
    replanning due to failures, delays, or environmental changes.
    
    \item To analyze the trade-offs between cloud-based and locally deployed models in terms of latency, reliability,
    and responsiveness during mission execution.
    
    \item To identify the limitations of LLM-based planning approaches and determine under which conditions they
    outperform or fall short of traditional planning methods.
\end{itemize}

\section{Thesis Structure}

\section{Thesis Structure}

This thesis is organized as follows:

Chapter 2 provides the technical background required for the thesis. It introduces autonomous
drone systems, multi-agent coordination and planning, vehicle routing concepts, Large Language
Models, and related work on LLM-based robotic and multi-agent planning.

Chapter 3 presents the proposed methodology and system architecture. It describes the
LLM-based planning pipeline and also explains the local task admission module, the replanning mechanism, and the software
architecture used to coordinate the planner and drone workers.

Chapter 4 presents the evaluation results, including cloud-based planner comparisons against a vehicle routing problem baseline, 
local task admission evaluation, and a simulation case study demonstrating planning, failure handling, and replanning. 
It also describes the physical Parrot Anafi platform and laboratory setup for future real-world deployment.

Chapter 5 concludes the thesis by summarizing the main findings, discussing the limitations
of the proposed approach, and outlining directions for future work.